\section*{Chapitre \ref{ch:b1:e-ph-diagrams} --- Diagrammes potentiel--pH}

\begin{solution}{exo:b1:e-ph-diagrams:1}
Manganèse~: \ce{Mn} (0) $<$ \ce{Mn^2+}, \ce{Mn(OH)2} ($+$II) $<$ \ce{MnO2}
($+$IV) $<$ \ce{MnO4-} ($+$VII). Chlore~: \ce{Cl-} ($-$I) $<$ \ce{Cl2} (0)
$<$ \ce{HClO}, \ce{ClO-} ($+$I) $<$ \ce{ClO3-} ($+$V).
\end{solution}

\begin{solution}{exo:b1:e-ph-diagrams:2}
$-0.059 \times 3/1 = -0.177$~; $+0.059 \times 2/2 = +0.059$~;
$-0.059 \times 8/5 = -0.094$~; $-0.059 \times 2/2 = -0.059$ V par unité
de pH. Pour le couple du cuivre, les protons sont libérés du côté de la
forme réduite ($m = -2$)~: le potentiel croît avec le pH.
\end{solution}

\begin{solution}{exo:b1:e-ph-diagrams:3}
\ce{H+}/\ce{H2} passe par 0~V à pH 0 et n'atteint jamais 0.50~V pour
$\mathrm{pH} \geqslant 0$. \ce{O2}/\ce{H2O} atteint 0.50~V à
$\mathrm{pH} = 0.73/0.059 = 12.4$ et 0~V seulement à pH 20.8, hors de
l'échelle. La bande a une largeur de \qty{1.23}{V} à tout pH, pH 7
compris (de $-0.41$ à $+0.82$~V).
\end{solution}

\begin{solution}{exo:b1:e-ph-diagrams:4}
(1, 1.0~V)~: \ce{Fe^3+}. (4, 0)~: la droite \ce{Fe(OH)3}/\ce{Fe^2+} est
à $1.09 - 0.71 = 0.38$~V et le point est au-dessous, au-dessus de
$-0.47$~V~:
\ce{Fe^2+}. (10, $-0.4$~V)~: entre $-0.07 - 0.59 = -0.66$ et $0.28 -
0.59 = -0.31$~V~: \ce{Fe(OH)2}. (10, 0.5~V)~: \ce{Fe(OH)3}.
\end{solution}

\begin{solution}{exo:b1:e-ph-diagrams:5}
\ce{Fe^2+}/\ce{Fe}~: $-0.41 + 0.030 \times (-6) = -0.59$~V.
\ce{Fe^3+}/\ce{Fe(OH)3}~: $14.00 - \frac13(38.55 - 6) = 3.15$. Les
domaines des ions dissous (corrosion) s'agrandissent~: \ce{Fe^2+}
descend jusqu'à $-0.59$~V et \ce{Fe^3+} s'étend jusqu'à pH 3.15~; le
métal et les hydroxydes perdent du terrain.
\end{solution}

\begin{solution}{exo:b1:e-ph-diagrams:6}
\ce{Fe(OH)3 + 3H+ + e- -> Fe^2+ + 3H2O}~: $E = E^\circ + 0.059\log\frac{[\ce{H+}]^3}{[\ce{Fe^2+}]}
= (E^\circ - 0.059\log c) - 0.177\,\mathrm{pH}$. À pH 1.82, la droite
rejoint \ce{Fe^3+}/\ce{Fe^2+} à 0.77~V, donc la constante vaut
$0.77 + 0.177 \times 1.82
= 1.09$~V.
\end{solution}

\begin{solution}{exo:b1:e-ph-diagrams:7}
$E^\circ(\ce{Cu+}/\ce{Cu}) = 0.52 > E^\circ(\ce{Cu^2+}/\ce{Cu+}) = 0.16$~:
d'après la \cref{prop:b1:e-ph-diagrams:disproportionation-criterion},
\ce{Cu+} se dismute. Frontière unique~: $E = 0.34 + 0.030\log 0.010 =
0.28$~V.
\end{solution}

\begin{solution}{exo:b1:e-ph-diagrams:8}
À pH 0, le domaine du cuivre (au-dessous de 0.28~V) chevauche celui de
l'eau et se trouve au-dessus de la droite du dihydrogène~: pas de
réaction avec l'acide chlorhydrique en l'absence d'air. Acide
nitrique~: \ce{NO3-}/\ce{NO}, à 0.96~V, est au-dessus du domaine du
cuivre, qui est oxydé~:
\ce{3Cu + 2NO3- + 8H+ -> 3Cu^2+ + 2NO + 4H2O}.
% ledger: dfg:NO3-_ao, dfg:NO_g (E0 computed in figdata/bachelor-1/standard-potentials.py)
\end{solution}

\begin{solution}{exo:b1:e-ph-diagrams:9}
À pH 3, la droite \ce{O2}/\ce{H2O} est à $1.23 - 0.18 = 1.05$~V,
au-dessus de tout le domaine de \ce{Fe^2+}~: le dioxygène oxyde le
fer(II). Au-dessus de pH 1.82, le fer(III) formé précipite sous forme
d'hydroxyde jaune-brun~:
\ce{4Fe^2+ + O2 + 10H2O -> 4Fe(OH)3 + 8H+}.
\end{solution}

\begin{solution}{exo:b1:e-ph-diagrams:10}
Verticales~: $14.00 - \frac12(16.38 - 2) = 6.81$~; et, à partir de
\ce{Zn(OH)2 + 2OH- <=> Zn(OH)4^2-} ($K = 10^{-1.93}$) avec
$[\ce{Zn(OH)4^2-}] = 10^{-2}$, $[\ce{OH-}]^2 = 10^{-0.07}$, pH 13.97.
\ce{Zn^2+}/\ce{Zn}~: $-0.76 + 0.030 \times (-2) = -0.82$~V.
\ce{Zn(OH)2}/\ce{Zn}~:
pente $-0.059$, passant par $(6.81, -0.82)$~: $E = -0.42 - 0.059\,\mathrm{pH}$.
\ce{Zn(OH)4^2- + 4H+ + 2e- -> Zn + 4H2O}~: pente $-0.118$, passant par
$(13.97, -1.24)$~: $E = 0.41 - 0.118\,\mathrm{pH}$.
\end{solution}

\begin{solution}{exo:b1:e-ph-diagrams:11}
Le domaine du fer est au-dessous de la droite du dihydrogène à tout pH.
pH 2~: le fer est oxydé en \ce{Fe^2+} avec dégagement de dihydrogène
(corrosion). pH 7~: la réaction reste possible, mais la droite du
dihydrogène ($-0.41$~V) est proche de \ce{Fe^2+}/\ce{Fe} et la réaction
est très lente. pH 13~: le fer est oxydé en \ce{Fe(OH)2}, \omterm{def:b1:e-ph-diagrams:corrosion-domains}{domaine de
passivation}~: une couche se forme et le clou garde son éclat. En
présence de dioxygène dissous, la droite \ce{O2}/\ce{H2O}, bien
au-dessus, oxyde le fer en fer(III) à tout pH~: la rouille, \ce{Fe(OH)3}
et ses oxydes déshydratés.
\end{solution}

\begin{solution}{exo:b1:e-ph-diagrams:12}
Le zinc a le domaine le plus bas~: au contact de l'acier, il est oxydé
le premier, \ce{2Zn + O2 + 2H2O -> 2Zn(OH)2} à pH 8, et les électrons
qu'il libère réduisent le dioxygène à la surface de l'acier, qui reste
métallique. Le magnésium ($E^\circ = -2.36$~V) est plus bas encore et
convient aussi~; le cuivre (0.34~V) est au-dessus du fer~: relié au
cuivre, c'est le fer qui serait oxydé, et plus vite que seul.
\end{solution}

\begin{solution}{pb:b1:e-ph-diagrams:1}
\textbf{1.} $-$I, 0, $+$I, $+$I.
\textbf{2.} \ce{Cl-} en bas, \ce{Cl2}, puis \ce{HClO} et \ce{ClO-} en
haut.
\textbf{3.} \ce{HClO} et \ce{ClO-}, un \omterm{def:b1:predominant-reaction:bronsted}{couple acide/base}.
\textbf{4.} \ce{Cl2 + 2e- -> 2Cl-}.
\textbf{5.} \ce{2HClO + 2H+ + 2e- -> Cl2 + 2H2O}.
\textbf{6.} \ce{ClO- + H2O + 2e- -> Cl- + 2OH-}.
\textbf{7.} \ce{HClO} au-dessous de pH 7.55, \ce{ClO-} au-dessus.
\textbf{8.} Aucun électron n'est échangé~; la frontière est fixée par la
seule $K_a$.
\textbf{9.} $1/(1 + 10^{7.4 - 7.55}) = 0.59$~: \qty{59}{\percent}.
\textbf{10.} $1/(1 + 10^{0.45}) = 0.26$~: les trois quarts du chlore
actif seraient sous forme de \ce{ClO-}, le désinfectant le moins
efficace.
\textbf{11.} \ce{ClO-}.
\textbf{12.} $E = 1.396 + 0.0296\log\frac{[\ce{Cl2}]}{[\ce{Cl-}]^2} = 1.396
+ 0.0296\log\frac{0.005}{10^{-4}} = 1.396 + 0.050 = 1.446$~V~; pas de
\ce{H+} dans la \omterm{def:b1:nernst:couple}{demi-équation}.
\textbf{13.} $E = 1.594 + 0.0296\log\frac{[\ce{HClO}]^2[\ce{H+}]^2}{[\ce{Cl2}]}
= 1.594 + 0.0296\log\frac{10^{-4}}{0.005} - 0.0592\,\mathrm{pH} = 1.594 -
0.050 - 0.0592\,\mathrm{pH}$.
\textbf{14.} En égalant, $0.0592\,\mathrm{pH} = 1.594 - 1.396 - 2 \times 0.0503
= 0.0974$, $\mathrm{pH} = 1.646$.
\textbf{15.} \ce{HClO + H+ + 2e- -> Cl- + H2O}~: un proton pour deux
électrons, pente $-0.0592/2 = -0.0296$.
\textbf{16.} En additionnant les deux \omterm{def:b1:nernst:couple}{demi-équations} (un électron par
atome de chlore dans chacune)~: $2E^\circ = 1.594 + 1.396$,
$E^\circ(\ce{HClO}/\ce{Cl-}) =
1.495$~V~; frontière $E = 1.495 - 0.0296\,\mathrm{pH}$.
\textbf{17.} \ce{ClO- + 2H+ + 2e- -> Cl- + H2O}~: $E = 1.718 - 0.0592\,\mathrm{pH}$
(constante choisie pour la continuité~: $1.495 - 0.0296 \times 7.55 = 1.272 = 1.718
- 0.0592 \times 7.55$).
\textbf{18.} \ce{Cl-} au-dessous des droites~; \ce{Cl2} dans un petit
triangle pour
$\mathrm{pH} < 1.65$ au-dessus de 1.446~V~; \ce{HClO} au-dessus de la
droite de pente
$-0.0296$ jusqu'à pH 7.55, \ce{ClO-} au-delà. La droite
\ce{O2}/\ce{H2O}, de 1.23 à 0.40~V, est au-dessous de toutes ces
frontières (le diagramme calculé par le script du chapitre a la même
allure).
\textbf{19.} Non~: son domaine est au-dessus de la droite du dioxygène,
il peut donc oxyder l'eau, mais la réaction est lente. Une piscine perd
son chlore surtout parce qu'il oxyde ce qu'apportent les baigneurs et
l'air, et sous l'effet du soleil~; il faut en rajouter.
\textbf{20.} \ce{Cl2} n'a pas de domaine dans cette zone~: il se dismute,
\ce{Cl2 + H2O -> HClO + Cl- + H+}~; en milieu basique,
\ce{Cl2 + 2OH- -> ClO- + Cl- + H2O}.
\textbf{21.} En faisant barboter du dichlore dans une solution
d'hydroxyde de sodium (la réaction de la question 20).
\textbf{22.} \ce{HClO + Cl- + H+ -> Cl2 + H2O}, une \omterm{def:b1:e-ph-diagrams:disproportionation}{médiamutation}.
\textbf{23.} Au-dessous de pH 1.6, l'hypochlorite et le chlorure
subissent une \omterm{def:b1:e-ph-diagrams:disproportionation}{médiamutation}~: du dichlore se forme et, au-delà de sa
\omterm{def:b1:precipitation:solubility}{solubilité}, se dégage sous forme de gaz toxique.
\textbf{24.} Pour $c = \qty{0.010}{mol/L}$, pH 4 est hors du domaine de
\ce{Cl2}. Mais la frontière dépend de $c$~: en refaisant les questions
12 à 14 pour un $c$ quelconque, on obtient
$\mathrm{pH} = 3.34 + \log(2c)$, qui atteint 3.6 pour
$c = \qty{1}{mol/L}$, l'ordre de grandeur de la concentration de l'eau
de Javel. Mélanger des produits concentrés vers pH 4 n'est pas sûr non
plus.
\textbf{25.} \textbf{pH 1.6}~: au-dessus, pour $c = \qty{0.010}{mol/L}$,
le dichlore dissous se dismute en acide hypochloreux et en chlorure.
\end{solution}
